\chapter{Synthesizer Basics}
\label{chap:synth-basics}

\classquote{不徐不疾，得之于手而应于心。}{《庄子·天道》 / Zhuangzi, "The Way of Heaven" (Tian Dao)}{Translation: Neither too fast nor too slow — the hand finds it and the mind follows. Note: synth parameters have no single correct value; balance what you hear with what you feel.}

This chapter explains the built-in poly synth's signal chain, its nine
projectable parameters, and the exact values of the six factory presets. The
parameter names, ranges, and defaults come from the engine's device projection
code; values are reproduced without any processing beyond rounding.

\section{Signal chain}

The built-in synth uses one fixed-size voice pool per track, holding at most 16
sounding voices. Each voice has two oscillator branches, a four-pole ladder
low-pass filter, and an amplitude envelope. Both branches draw from the same five
wavetables: \texttt{hollow}, \texttt{organ}, \texttt{vocal},
\texttt{metallic}, and \texttt{glass}. Each branch has its own level and detune;
the branches are mixed before the low-pass and envelope.

\begin{center}
\resizebox{\textwidth}{!}{\begin{tikzpicture}[
  node distance=9mm and 11mm,
  box/.style={draw, rounded corners=2pt, minimum width=24mm, minimum height=8mm, align=center, font=\small},
  sum/.style={draw, circle, minimum size=6mm, inner sep=0pt, font=\small},
  arrow/.style={-{Latex[length=2mm]}, thick}
]
  \node[box] (osc1) {OSC 1\\wavetable + level + detune};
  \node[box, below=of osc1] (osc2) {OSC 2\\wavetable + level + detune};
  \node[sum, right=of osc1, yshift=-8mm] (mix) {+};
  \node[box, right=of mix] (filter) {Four-pole\\ladder low-pass};
  \node[box, right=of filter] (env) {Amplitude envelope\\fixed ADSR};
  \node[box, right=of env] (out) {Insert chain\\and bus};

  \draw[arrow] (osc1) -- (mix);
  \draw[arrow] (osc2) -- (mix);
  \draw[arrow] (mix) -- (filter);
  \draw[arrow] (filter) -- (env);
  \draw[arrow] (env) -- (out);
\end{tikzpicture}}
\end{center}

The filter is a four-pole ladder low-pass. The device clamps cutoff to between
20 Hz and 0.45 times the sample rate; resonance and drive are 0-to-1 knob values.
At drive 0 the distortion stage is bit-exactly bypassed; at 1 it is +36 dB of
pre-gain, distributed evenly in dB in between. The output is self-compensated, so
turning up drive changes harmonics rather than overall loudness.

The amplitude envelope is a fixed ADSR: 5 ms attack, 80 ms decay, 0.7 sustain,
and 50 ms release. The current interface exposes no adjustable ADSR knobs, and no
LFO, modulation matrix, or separate envelope-stage controls. Do not apply control
names from other synthesizers here. If these are wired later, find the
corresponding parameter projection in the code first, then extend this chapter.

\section{The nine projectable parameters}

The built-in synth reads device-chain parameter names case-insensitively. The
table lists the nine parameters, their valid ranges, and defaults. ``Absent'' and
``written as 0'' are different: an absent \texttt{osc2\_level} means the second
branch is off, and an explicit 0 also means off; an absent \texttt{cutoff\_hz}
bypasses the whole filter, while an explicit 0 is clamped to 20 Hz by the device.

\begin{table}[htbp]
\centering
\small
\begin{tabular}{@{}llll@{}}
\hline
Parameter & Meaning & Valid range & Default when absent \\
\hline
\texttt{cutoff\_hz} (alias \texttt{cutoff}) & Low-pass cutoff & 20 Hz--0.45×sample rate & Filter bypassed \\
\texttt{resonance} (alias \texttt{res}) & Low-pass resonance & 0--1 & 0.0 \\
\texttt{drive} & Filter drive/distortion & 0--1 & 0.0 \\
\texttt{osc1\_wave} & First-branch wavetable & 0--4 (rounded) & 0 (hollow) \\
\texttt{osc2\_wave} & Second-branch wavetable & 0--4 (rounded) & 0 (hollow) \\
\texttt{osc1\_level} & First-branch level & 0--1 & 1.0 (full) \\
\texttt{osc1\_detune\_cents} & First-branch detune & ±4800 cents & 0.0 \\
\texttt{osc2\_level} & Second-branch level & 0--1 & 0.0 (off) \\
\texttt{osc2\_detune\_cents} & Second-branch detune & ±4800 cents & 0.0 \\
\hline
\end{tabular}
\caption{The nine parameters of the built-in poly synth}
\label{tab:synth-params}
\end{table}

\screenshot{synth-device-params}{Synth nine parameters and device-card readout}

The wavetable index mapping is fixed: 0 = \texttt{hollow}, 1 = \texttt{organ},
2 = \texttt{vocal}, 3 = \texttt{metallic}, 4 = \texttt{glass}. Non-finite or
negative values are treated as 0; integers above 4 are clamped by the device to
the last table, not turned into silence. When a branch level is 0, that branch
executes no samples at all; this is a bit-exact bypass, not a multiply by zero.

\section{Parameter values of the six factory presets}

The two tables below give all nine parameters for each preset. There is no
``envelope'' or ``LFO'' column because the current presets contain no such
parameters.

\begin{table}[htbp]
\centering
\small
\begin{tabular}{@{}lrrr@{}}
\hline
Preset & \texttt{cutoff\_hz} & \texttt{resonance} & \texttt{drive} \\
\hline
Hollow Lead & 2400 & 0.20 & 0.00 \\
Warm Organ & 1800 & 0.15 & 0.10 \\
Choir Pad & 2200 & 0.25 & 0.00 \\
FM Bell & 4200 & 0.30 & 0.15 \\
Glass Pluck & 3600 & 0.20 & 0.00 \\
Sub Bass & 420 & 0.35 & 0.25 \\
\hline
\end{tabular}
\caption{Filter parameters of the factory presets}
\label{tab:preset-filter}
\end{table}

\begin{table}[htbp]
\centering
\small
\begin{tabular}{@{}lrrrrrr@{}}
\hline
Preset & \texttt{osc1\_wave} & \texttt{osc2\_wave} & \texttt{osc1\_level} & \texttt{osc1\_detune} & \texttt{osc2\_level} & \texttt{osc2\_detune} \\
\hline
Hollow Lead & 0 & 0 & 0.90 & $-4$ & 0.35 & $+5$ \\
Warm Organ & 1 & 1 & 0.80 & $-5$ & 0.60 & $+6$ \\
Choir Pad & 2 & 0 & 0.85 & $-8$ & 0.55 & $+8$ \\
FM Bell & 3 & 4 & 0.85 & $-3$ & 0.50 & $+11$ \\
Glass Pluck & 4 & 4 & 0.85 & $-4$ & 0.45 & $+4$ \\
Sub Bass & 0 & 1 & 0.95 & $0$ & 0.35 & $-5$ \\
\hline
\end{tabular}
\caption{Oscillator parameters of the factory presets; detune is in cents}
\label{tab:preset-osc}
\end{table}

Tables \ref{tab:preset-filter} and \ref{tab:preset-osc} give each preset its
character directly: Hollow Lead is a 90\%-level hollow branch plus a low-level
second branch of the same table, suited to a soft lead; Warm Organ opens both
organ branches and uses detune for organ-like thickness; Choir Pad uses vocal as
the main branch and hollow as the auxiliary branch, and its pad character comes
from the vocal table combined with a lower cutoff; FM Bell and Glass Pluck both
use metallic/glass tables at high cutoff, the former brighter and the latter
shorter; Sub Bass pushes cutoff down to 420 Hz and keeps only the low end.

\subsection{Three landmark synths: three characters of subtractive synthesis}

A parameter table tells you what is there now; it does not tell you why to play it that way. Three instruments in the history of subtractive synthesis map onto this chapter's three moves: the Minimoog turns ``less'' into weight, the Juno turns ``thin'' into width, and the TB-303 turns ``one line'' into groove.

\textbf{The Minimoog (1970)} packed the room-sized modular Moog into a portable instrument: three oscillators, a ladder low-pass, and no recallable presets. Its lesson is not ``more parameters are better'', but ``every branch needs a job''. This chapter has only two branches, which is exactly the exercise: let \texttt{osc1} establish the line, then use \texttt{osc2}'s level and detune for thickness; leave the third layer to arrangement instead of forcing it into one patch. Yeban's Sub Bass, with a 420 Hz low-pass and one organ branch, is another way of writing ``less''.

\textbf{The Juno-60/Juno-106 (1982/1984)} used DCOs and chorus to turn ``thin'' into ``wide''. Its chorus does not add a real third oscillator; it modulates time and stereo placement. This chapter has no chorus effect, so you can only use small detune between two branches of the same table, say $\pm5$ to $\pm9$ cents, and leave the rest of the width to later effects. Do not claim this equals the Juno chorus; it is an approximation with the available parameters, not a replica.

\textbf{The TB-303 (1981)} was designed as a single-oscillator-plus-filter bass instrument for bass players, but the dance floor turned it into an acid lead. It reminds us that resonance and cutoff are not decoration: they change pitch perception and groove. Push \texttt{resonance} from 0.2 to 0.5 and the harmonics around cutoff are emphasised, so the melodic centre of gravity moves; that is not volume, and \texttt{drive} should not be used to compensate. Yeban has no 303 glide or accent controls; what you can practise is the judgment of ``one line plus one filter''.

The common lesson: make one line clear first, then add the second; listen to phase and detune before talking about thickness. Parameters are hands; listening is the ear. Do not let a preset name listen for you.

\section{Help-topic mapping for this chapter}

\begin{table}[htbp]
\centering
\small
\begin{tabular}{@{}p{0.34\linewidth}p{0.52\linewidth}@{}}
\hline
Section & In-app help anchor \\
\hline
Signal chain & \texttt{synth-signal-chain} \\
Nine projectable parameters & \texttt{synth-parameters} \\
Parameter values of the six factory presets & \texttt{synth-factory-presets} \\
\hline
\end{tabular}
\caption{Help-topic mapping for Chapter 2}
\end{table}

\section*{Further Reading}
\begin{itemize}[leftmargin=1.5em]
\item \url{https://www.soundonsound.com/techniques} — Sound On Sound techniques: compare oscillators, filters, and envelopes against working tutorials and real listening examples.
\item \url{https://www.ableton.com/en/blog/} — Ableton blog: read synth patch breakdowns and producer workflows, and fill in the musical judgment behind the parameters.
\end{itemize}
